\medskip

\begin{enumerate}
\item On considère la suite $(u_n)$ définie pour tout $n$ non nul par 
$u_n = \dfrac{25 + (- 1)^n}{n}.$

\textbf{Affirmation 1 :} La suite $(u_n)$ est divergente.

\medskip

Pour tout $n$ non nul, on a:
$ -1 \leqslant (-1)^n \leqslant +1$
donc $ 25-1 \leqslant 25+(-1)^n \leqslant 25+1$
donc $ 24 \leqslant 25+(-1)^n \leqslant 26$
donc $ \dfrac{24}{n} \leqslant \dfrac{25+(-1)^n}{n} \leqslant \dfrac{26}{n}$
donc $ \dfrac{24}{n} \leqslant u_n \leqslant \dfrac{26}{n}$

Or $\ds\lim_{n\to +\infty} \dfrac{24}{n} = \ds\lim_{n\to +\infty} \dfrac{26}{n} = 0$ donc, d'après le théorème des gendarmes, on peut déduire que la suite $(u_n)$ est convergente et que $\ds\lim_{n\to +\infty} u_n=0$.

\hfill \textbf{Affirmation 1 fausse}

\medskip

\item On considère la suite $(w_n)$ définie pour tout entier naturel $n$ par $\left\{\begin{array}{l !{=}  l}
w_0&1\\w_{n+1}&\dfrac{w_n}{1 + w_n}
\end{array}\right.$

On admet que pour tout entier naturel $n,\; w_n > 0$n et on considère la suite $(t_n)$ définie pour tout entier naturel $n$ par $t_n =\dfrac{k}{w_n}$ où $k$ est un nombre réel strictement positif.

\textbf{Affirmation 2 :} La suite $(t_n)$ est une suite arithmétique strictement croissante.

\medskip

Pour tout $n$, on a: $t_{n+1}
= \dfrac{k}{w_{n+1}}
= \dfrac{k}{\frac{w_n}{1+w_n}}
= \dfrac{k\left (1+w_n\right )}{w_n}
= \dfrac{k+kw_n}{w_n}
=\dfrac{k}{w_n} +k
=t_n+k$

Donc la suite $(t_n)$ est arithmétique de raison $k$.

De plus, $k>0$ donc la suite $(t_n)$ est strictement croissante.

\hfill \textbf{Affirmation 2 vraie}

\medskip

\item On considère la suite $(v_n)$ définie pour tout entier naturel $n$ par 
$\left\{\begin{array}{l !{=} l}
v_0&1\\
v_{n+1}&\ln \left(1 +  v_n\right)
\end{array}\right.$

On admet que pour tout entier naturel $n$, $v_n > 0$,

\textbf{Affirmation 3 :} La suite $(v_n)$ est décroissante.

\medskip

Soit $\mathcal{P}_n$ la propriété: $v_n>v_{n+1}$.

On va démontrer par récurrence que cette propriété est vraie pour tout $n$.

\begin{list}{\textbullet}{}
\item \textbf{Initialisation}

$v_0=1$ et $v_1=\ln\left (1+v_0\right ) = \ln\left (2\right ) \approx 0,69$; donc $v_0>v_1$.\\
La propriété est vraie pour $n=0$.

\item \textbf{Hérédité}

On suppose que $v_n > v_{n+1}$; c'est l'hypothèse de récurrence.

$v_n > v_{n+1}$ donc $1+v_n > 1+v_{n+1}$

Or la fonction ln est  croissante sur $]0\,;\, +\infty[$ donc
$\ln\left (1+v_n\right ) >\ln\left ( 1+v_{n+1}\right )$, ce qui veut dire que
$v_{n+1} > v_{n+2}$.
La propriété est donc héréditaire.

\item \textbf{Conclusion}

La propriété est vraie au rang 0. Elle est héréditaire pour tout $n\geqslant 0$. Donc, d'après le principe de récurrence, elle est vraie pour tout $n\geqslant 0$.
\end{list}

On a démontré que, pour tout $n\geqslant 0$, on avait: $v_{n} > v_{n+1}$; donc la suite $(v_n)$ est décroissante.

\hfill \textbf{Affirmation 3 vraie}

\medskip

\item On considère la suite $(I_n)$ définie pour tout entier naturel $n$ par $I_n = \displaystyle\int_1^{\e} \left [\ln (x)\strut\right ]^n\,\d x$.

\textbf{Affirmation 4 :} Pour tout entier naturel $n$, $I_{n+1} = \e - (n + 1) I_n$.

\medskip

$I_{n+1} = \displaystyle\int_1^{\e} \left [\ln (x)\strut \right ]^{n+1}\d x
= \displaystyle\int_1^{\e} 1\times \left [\ln (x)\strut\right ]^{n+1}\d x$

On va calculer $I_{n+1}$ au moyen d'une intégration par parties:\\
$\ds\int_{a}^{b} u'(x)v(x) \d x = \left [u(x) v(x)\strut \right ]_{a}^{b} - \ds\int_{a}^{b} u(x) v'(x) \d x$

On pose $u'(x)=1$ donc $u(x)=x$, et $v(x)= \left [\ln\left (x\right )\strut\right] ^{n+1}$ donc
$v'(x)=\left (n+1\right ) \times \dfrac{1}{x} \times \left [\ln\left (x\right )\strut\right ]^n$.

$\begin{aligned}[t]
I_{n+1} & =  \left [ x\times \left [\ln\left (x\right )\strut\right ]^{n+1}\right ]_{1}^{\e}
 - \ds\int_{1}^{\e} x\times (n+1) \times \dfrac{1}{x} \times \left [ \ln\left (x\right )\strut \right ]^n \d x\\
 & = \left ( \e \left [ \ln \left (\e\right )\strut \right ]^{n+1} - 1 \left [ \ln \left (1\right )\strut \right ]^{n+1}\right ) - \left (n+1\right ) \ds\int_{1}^{\e} \left [\ln\left (x\right )\strut\right ]^{n} \d x\\
&  = \e - \left (n+1\right ) I_n
 \end{aligned}$

\hfill \textbf{Affirmation 4 vraie}
\end{enumerate}

\vspace{0,5cm}
